\section{Behavioural Repository Sufficiency}
\label{sec:sufficiency}

The original v0.1 draft expressed repository sufficiency through equality of latent distributions over “useful next actions” and a KL-divergence bound. Hostile review rejects that formulation for the planned experiments. The action space is open-ended, equivalent implementations need not have comparable surface actions, and the proposed runs do not expose a well-defined probability distribution over engineering actions. KL divergence would therefore add mathematical appearance without an estimable quantity.

The paper instead defines an observable behavioural target.

Let $Y_{s,d,c}\in\{0,1\}$ denote preregistered hidden-verifier success for task sequence $s$, depth $d$, and condition $c$. Let $c=A$ be the persistent-history control and $c=B$ the forced-reset condition. The two conditions must use the same accepted starting repository state, current task, high-capability model/configuration, tool contract, workspace image, and acceptance criteria. The intended treatment difference is availability of predecessor reasoning-session history.

For later repeated confirmatory experiments, define:
\begin{equation}
p_c(d)=P(Y=1\mid c,d,\mathcal{M},\mathcal{E},\mathcal{T}),
\label{eq:behavioural-success}
\end{equation}
where $\mathcal{M}$ identifies the matched model/configuration, $\mathcal{E}$ the controlled execution environment, and $\mathcal{T}$ the task-sampling population.

A future confirmatory study can preregister a scientifically justified non-inferiority margin $\delta>0$ and ask whether:
\begin{equation}
p_B(d)-p_A(d) \ge -\delta
\label{eq:noninferiority-target}
\end{equation}
over the declared task population and depth range. If symmetric equivalence is scientifically required, a two-sided margin can be preregistered instead. The margin must be fixed before outcomes and justified in engineering terms; this paper does not invent one.

P0 is too small to establish non-inferiority statistically. Its role is to test whether the intervention, task corpus, hidden verification, leakage controls, and telemetry are workable and whether an obvious continuity collapse appears.

“Repository sufficiency” is therefore shorthand for a conditional behavioural claim, not an intrinsic property of a repository. A repository may be adequate for one task class, model, tool interface, and depth but not another. It also need not preserve every past thought. The question is whether accepted durable state permits matched fresh reasoning to achieve comparable engineering outcomes at acceptable reconstruction cost.

Operational evidence includes:
\begin{itemize}
    \item hidden-verifier task success;
    \item complete-sequence success across dependent tasks;
    \item success by sequential depth;
    \item reconstruction-probe accuracy;
    \item regression rate;
    \item human-independent resumption;
    \item reconstruction resource use and elapsed time.
\end{itemize}

Semantic-state ablation remains useful but must test marginal continuity value rather than “making the repository worse”. Condition C should remove one small preregistered artefact class or information channel at a time while keeping task semantics and buildability intact. Redundant encoding must be audited: deleting an ADR proves little if the same decision is reproduced verbatim in a test, comment, or later allowed commit.

The strong RaS hypothesis is thus intentionally narrow: for a declared class of dependent software-engineering tasks at accepted transition boundaries, predecessor reasoning-session history may have sufficiently small marginal behavioural value that a fresh reasoner can continue from authoritative project state without material outcome loss. Whether that is true is an empirical question.

\subsection{What the accepted Subject-B evidence shows}

Subject-B P1 observed equal non-zero behavioural performance: 18/30 behaviours in each condition, with 30 matched agreements and no disagreements across three tasks. Subject-B P2 expanded the same-repository design to four tasks and three repetitions. It observed 0/39 satisfied behaviours in both conditions, 39 matched agreements, no disagreements, and 0/12 candidate-level overall passes in each condition. The P2 result is therefore a floor effect: it detects no A/B difference under this design but does not demonstrate successful behavioural preservation or repository-state sufficiency. Neither result supports an equivalence or non-inferiority claim, and the two studies must not be pooled into one correctness score.
